
\chapter{Conclusion and Future Work}\label{ch:conclusion}

\section{Summary}

This thesis presented a security-first fork of the OCI reference runtime \texttt{runc}. The fork makes three concrete changes to the runtime layer: it ships with an empty default capability set (with two opt-in flags that restore the upstream and the historical Docker defaults), it adds a single \texttt{-{}-security-scan} flag that records the workload through eBPF and emits seccomp, AppArmor, and capability profiles, and it narrows the on-disk \texttt{config.json} to the observed capability set without ever widening it.

Two design decisions distinguish the implementation from prior automated-profiling work. First, every hook is a hidden sub-command of the runtime binary itself, so the OCI bundle never gains an executable artefact at scan time. Second, capability tracing uses \texttt{capable-bpfcc -{}-cgroupmap} with a \texttt{bpftool}-pinned cgroup-v2 map, so child processes spawned after the trace started are observed --- the failure mode that simpler \texttt{-p \$pid} pipelines silently exhibit.

The proposed runtime was situated against manual policy authoring, cluster-side recorders such as the Kubernetes Security Profiles Operator and Inspektor Gadget, and sandboxed runtimes (gVisor, Kata) analysed from the literature. The argument is not that the runtime dominates every alternative: a skilled operator or a staffed DevSecOps pipeline can in principle exceed its output. The argument is that the runtime substantially lowers the operational floor for teams that run plain \texttt{runc} without orchestration.

\section{Answers to the Research Questions}

\textbf{RQ1 --- Can a runtime, without infrastructure changes, generate profiles whose security level approaches manual tuning or cluster-side recorders?} On exercised paths, yes: the laboratory produced deny-by-default seccomp, AppArmor rules from complain-mode audit records, and cgroup-wide capability traces from a single \texttt{--security-scan} run (Chapter~\ref{ch:experiments}, Table~\ref{tab:profquality-lab}). File capabilities on executed binaries (e.g.\ \texttt{CAP\_NET\_RAW} on \texttt{ping}) are not yet merged from \texttt{security.capability} xattrs; that limitation is documented as a drawback rather than hidden.

\textbf{RQ2 --- Does scan-then-enforce prevent command-injection exploitation and improve containment relative to stock \texttt{runc}?} On the tested workload, scan-then-enforce prevented command injection and blocked reconnaissance, sensitive reads, reverse shells, mount probes, and persistence attempts that succeeded under baseline stock \texttt{runc}, while preserving legitimate ping and calendar behaviour after documented \texttt{apply.sh} adjustments (Chapter~\ref{ch:experiments}). The application flaw remains in the image; the runtime stops its exploitation at the syscall boundary.

\section{Limitations}

\textbf{Dynamic-trace coverage.} Like every behavioural profiler --- including SPO and Inspektor Gadget --- the scanner only records what executes during the scan window. Code paths gated by uncommon errors, scheduled jobs, signal handlers, log-rotation events, certificate refresh, database migrations, or fallback branches will not be recorded unless explicitly exercised. The recommended mitigation is to drive the scan with the same end-to-end test suite that gates the application's releases.

\textbf{Root workloads.} \texttt{cap\_capable} is not invoked by the kernel for many checks performed by \texttt{root}, so a workload that runs as \texttt{uid 0} systematically under-reports the capability set. The runtime warns but does not rewrite the user. Workloads that genuinely need root for legitimate reasons can either re-run the scan as a non-zero UID or accept the under-reporting as a known limitation.

\textbf{Cgroup-v2 only.} The cgroup-wide capability tracer requires the unified cgroup-v2 hierarchy and a recent \texttt{capable-bpfcc} that accepts \texttt{-{}-cgroupmap}, plus \texttt{bpftool}. Hosts on cgroup-v1 or hybrid hierarchies are explicitly refused; a silent fallback to \texttt{-p \$pid} would have hidden the regression in child-process coverage and was therefore rejected.

\textbf{Supply-chain malicious behaviour.} An image whose behaviour is intentionally crafted to look benign during the scan and only attempts privileged operations later is out of scope. This is a generic limitation of trace-based profiling and is shared with every cluster-side recorder.

\textbf{AppArmor refinement.} The runtime emits an AppArmor profile from complain-mode audit records and flips it into enforce mode when rules were learned. Operators may still review or further refine that profile with host tools such as \texttt{aa-logprof}.

\section{Future Work}

\textbf{Performance evaluation.} A dedicated benchmark harness (\texttt{performance-evaluation}) should compare stock \texttt{runc}, the hardened binary, enforcement on a scanned bundle, and sandboxed runtimes (gVisor, Kata) on startup latency, scalability, and sysbench workloads --- extending the analytical comparison in this thesis with empirical numbers on the same host.

\textbf{Broader laboratory coverage.} Additional intentionally vulnerable workloads, command-execution patterns, and attack scripts would test whether containment results generalise beyond the Flask system-monitor bundle.

\textbf{Quantitative containment metrics.} Future experiments should complement the current qualitative succeed/block outcomes with normalised metrics across attack classes and deployment mistakes, including privileged containers and \texttt{docker.sock} mounts when those scenarios are intentionally in scope.

\textbf{File-capability merge.} Reading \texttt{security.capability} xattrs on executed binaries during scan would close the \texttt{ping}/\texttt{CAP\_NET\_RAW} gap without laboratory workarounds.

\textbf{Static augmentation of the trace.} Combining the dynamic trace with static analysis of the image (entry points, ELF symbol tables, embedded scripts) would catch some of the missed code paths. Canella et al.~\cite{canella2021} already demonstrated this for seccomp; extending it to capabilities and AppArmor is a natural follow-up.

\textbf{Differential profiles across releases.} Re-running the scan on every image release and reporting the diff (new syscalls, new capabilities) would catch regressions before they ship.

\textbf{SPO compatibility.} Emitting the profiles in the SPO Custom Resource format would make them directly consumable by clusters that already run SPO.

\section{Closing Remarks}

The container-security literature is dominated by two endpoints: replace the runtime with a stronger sandbox, or wrap the runtime in a cluster-side platform. Both are sound; both also assume that someone is paying for the overhead --- in CPU for the sandbox, in engineer-hours for the platform. This thesis argues that the runtime layer itself is a viable third place to put confinement, and that doing so brings developer-grade security automation within reach of teams that cannot afford either of the existing options.
